\documentclass[../../../include/open-logic-section]{subfiles}

\begin{document}

\olfileid{sth}{ord-arithmetic}{using-addition}
\olsection{د ترتيبي جمع کارول}

د ترتيبي عددونو په جمع سره د بېلابېلو سټونو رتبې په ښکاره حسابولے شو،
د \olref[spine][rank]{defnsetrank} په معنا:

\begin{lem}\ollabel{rankcomputation}
که $\setrank{A} = \alpha$ او $\setrank{B} = \beta$ وي، نو:
\begin{enumerate}
	\item\ollabel{exrankpow} $\setrank{\Pow{A}} = \alpha\ordplus 1$
	\item\ollabel{exrankpair} $\setrank{\{A, B\}} = \max(\alpha,
	\beta) \ordplus 1$
	\item\ollabel{exrankcup} $\setrank{A \cup B} = \max(\alpha,
	\beta)$
	\item\ollabel{exranktuple} $\setrank{\tuple{A,B}} = \max(\alpha,
	\beta) \ordplus  2$
	\item\ollabel{exranktimes} $\setrank{A \times B} \leq \max(\alpha,
	\beta) \ordplus  2$
	\item\ollabel{exrankunion} $\setrank{\bigcup A} = \alpha$ چې
	$\alpha$ تش يا حدي وي؛ $\setrank{\bigcup A} = \gamma$ چې
	$\alpha = \gamma\ordplus 1$ وي
\end{enumerate}
\end{lem}

\begin{proof}
په ټول ثبوت کښې \olref[spine][rank]{ranksupstrict} بيا بيا کاروو۔

\emph{\olref{exrankpow}.} که $x \subseteq A$ وي، نو $\setrank{x} \leq
\setrank{A}$۔ نو $\setrank{\Pow{A}} \leq \alpha \ordplus  1$۔
څرنګه چې په ځانګړې توګه $A \in \Pow{A}$، نو
$\setrank{\Pow{A}} = \alpha \ordplus  1$۔

\emph{\olref{exrankpair}.} د \olref[spine][rank]{ranksupstrict} له مخې۔

\emph{\olref{exrankcup}.} د \olref[spine][rank]{ranksupstrict} له مخې۔

\emph{\olref{exranktuple}.} \olref{exrankpair} دوه ځله وکاروئ۔

\emph{\olref{exranktimes}.} پام وکړئ چې
$A \times B \subseteq \Pow{\Pow{A \cup B}}$، او \olref{exranktuple}
وکاروئ۔

\emph{\olref{exrankunion}.} که $\alpha = \gamma\ordplus 1$ وي، نو
داسې $c \in A$ شته چې $\setrank{c} = \gamma$ وي، او د $A$
هېڅ !!{element} تر دې لوړه رتبه نه لري؛ نو
$\setrank{\bigcup A} = \gamma$۔ که $\alpha$ حدي ترتيبي عدد وي،
نو $A$ داسې !!{element}s لري چې رتبې يې $\alpha$ ته هر څومره نږدې،
خو په کلکه ترې ټيټې، رسېږي؛ نو $\bigcup A$ هم داسې !!{element}s
لري چې رتبې يې $\alpha$ ته هر څومره نږدې، خو په کلکه ترې ټيټې،
رسېږي؛ له دې امله $\setrank{\bigcup A} = \alpha$۔
\end{proof}
\noindent
دا د تمرين دپاره پرېږدو چې ولې \olref{exranktimes} يو
\emph{نا}مساوات دے۔

\begin{prob}
داسې سټونه $A$ او $B$ جوړ کړئ چې $\setrank{A \times B}=
\max(\setrank{A}, \setrank{B})$ وي۔ داسې سټونه $A$ او $B$
جوړ کړئ چې
$\setrank{A \times B}=\max(\setrank{A}, \setrank{B}) \ordplus  2$
وي۔ ايا نورې رتبې هم ممکنې دي؟
OLSTH-015: د سرچينې د تمرين په دويمه معادله کښې د دواړو خواوو
تر منځ نښه نه وه؛ د مخکنۍ لمې د پورته کرانې د مثال دپاره دلته
مساوات په ښکاره ور زيات شوے دے۔ د تمرين نور شرطونه نه دي بدل شوي۔
\end{prob}

اوس دا هم ښودلے شو چې د يو ترتيبي عدد د «متناهي» يا «نامتناهي»
بللو څو معقولې معناوې سره برابرې دي:

\begin{lem}\ollabel{ordinfinitycharacter}
د هر ترتيبي عدد $\alpha$ دپاره لاندې شرطونه همارزښته دي:
\begin{enumerate}
	\item\ollabel{ord:notinomega} $\alpha\notin \omega$، يعنې
	$\alpha$ طبيعي عدد نه دے
	\item\ollabel{ord:omegaplus} $\omega \leq \alpha$
	\item\ollabel{ord:oneplus} $1 \ordplus  \alpha = \alpha$
	%\alpha \approx \alpha \ordplus 1$
	\item\ollabel{ord:plusone} $\alpha \approx \alpha\ordplus 1$،
	يعنې $\alpha$ او $\alpha\ordplus 1$ هم‌شمېره دي
	\item\ollabel{ord:infinite} $\alpha$ د ډېډېکېنډ په معنا نامتناهي دے
	\end{enumerate}
\end{lem}
\noindent
نو د ترتيبي عدد د (نا)متناهي والي د پوهېدو پنځه په ثبوت برابرې
لارې لرو۔

\begin{proof}
\emph{\olref{ord:notinomega} $\Rightarrow$ \olref{ord:omegaplus}.}
د درې‌ګوني وېش له مخې۔

\emph{\olref{ord:omegaplus} $\Rightarrow$ \olref{ord:oneplus}.}
$\alpha \geq \omega$ وټاکئ۔ د نامتناهيت هاخوا استقرا له مخې داسې
تر ټولو کوچنے ترتيبي عدد $\gamma$ (ښايي $0$) شته چې يو حدي
ترتيبي عدد $\beta$ د $\alpha = \beta \ordplus \gamma$ شرط پوره
کوي۔ اوس:
\[
	1 \ordplus \alpha =
	1 \ordplus (\beta \ordplus \gamma) =
	(1 \ordplus \beta) \ordplus \gamma =
	\supstrict_{\delta < \beta} (1 \ordplus  \delta) \ordplus  \gamma =
	\beta \ordplus  \gamma =
	\alpha.
\]
\emph{\olref{ord:oneplus} $\Rightarrow$ \olref{ord:plusone}.}
په څرګنده يوه !!a{bijection}
$f \colon (\alpha \disjointsum 1) \to (1 \disjointsum \alpha)$
شته۔ که $1 \ordplus \alpha = \alpha$ وي، نو يوه هم‌شکلتيا
$g \colon (1 \disjointsum \alpha) \to \alpha$ شته۔ اوس
$\comp{f}{g}$ وګورئ۔

\emph{\olref{ord:plusone} $\Rightarrow$ \olref{ord:infinite}.}
که $\alpha \approx \alpha \ordplus 1$ وي، يوه !!a{bijection}
$f \colon (\alpha \disjointsum 1) \to \alpha$ شته۔
$g(\gamma) = f(\gamma, 0)$ د هر $\gamma < \alpha$ دپاره تعريف کړئ؛
دا !!{injection} ښيي چې $\alpha$ د ډېډېکېنډ په معنا نامتناهي دے،
ځکه $f(0,1) \in \alpha \setminus \ran{g}$۔

\emph{\olref{ord:infinite} $\Rightarrow$ \olref{ord:notinomega}.}
دا \olref[z][infinity-again]{naturalnumbersarentinfinite} دے۔
\end{proof}

\end{document}
